% Manovella dello schienale (quesito ud2-rigido-manovella): r va dal perno O
% all'impugnatura, lungo la manovella; la forza F forma 30 gradi con la
% manovella. Il braccio è la distanza fra O e la retta d'azione di F
% (tratteggiata), b = r sin 30.
\begin{tikzpicture}[>=Stealth, font=\small]
  \coordinate (O) at (0,0);
  \coordinate (P) at (4,0);
  % manovella
  \draw[richiamo!40, line width=6pt, line cap=round] (O) -- (P);
  \draw[->, formula, very thick] (O) -- (P) node[midway, above=3pt] {$\vec r$};
  \fill[richiamo!70] (O) circle (0.16);
  \fill (O) circle (1.4pt);
  \node[left=3pt] at (O) {$O$};
  % forza e sua retta d'azione
  \draw[dashed, richiamo] ($(P)+(30:-4.0)$) -- (P);
  \draw[->, attenzione, very thick] (P) -- ++(30:1.6) node[above] {$\vec F$};
  \draw[dotted] (P) -- ++(1.3,0);
  \draw ($(P)+(0.9,0)$) arc (0:30:0.9);
  \node at ($(P)+(15:1.25)$) {\ang{30}};
  % braccio, perpendicolare alla retta d'azione
  \coordinate (H) at ($(P)+(30:-3.464)$);
  \draw[formula, thick] (O) -- (H) node[midway, left] {$b = r\sin\ang{30}$};
  \draw[formula] ($(H)+(120:0.22)$) -- ++(30:0.22) -- ($(H)+(30:0.22)$);
\end{tikzpicture}
